\newpage
\phantomsection
\label{prf:closed-set-characterization-topologies}
\LRAProofFor{thm:closed-set-characterization-topologies}

\begin{remark*}[Return]
\hyperref[thm:closed-set-characterization-topologies]{Return to Theorem}
\end{remark*}

\begin{theorem*}[Closed-Set Characterization of Topologies]
Let \(X\) be a set, and let \(\mathcal{C}\) be a collection of subsets of
\(X\).  Suppose that \(\emptyset,X\in\mathcal{C}\), arbitrary intersections
of elements of \(\mathcal{C}\) belong to \(\mathcal{C}\), and unions of two
elements of \(\mathcal{C}\) belong to \(\mathcal{C}\).  Then there is a
unique topology \(\tau\) on \(X\) whose family of closed sets is exactly
\(\mathcal{C}\).
\end{theorem*}

\begin{proof}[Professional Standard Proof]
\LRAProofBodyStart
Define
\[
  \tau=\{X\setminus C:C\in\mathcal{C}\}.
\]
Since \(X\in\mathcal{C}\) and \(\emptyset\in\mathcal{C}\), the sets
\(\emptyset=X\setminus X\) and \(X=X\setminus\emptyset\) belong to \(\tau\).

Let \(\{U_i\}_{i\in I}\) be any family of elements of \(\tau\).  For each
\(i\in I\), choose \(C_i\in\mathcal{C}\) such that \(U_i=X\setminus C_i\).
Then De Morgan's law gives
\[
  \bigcup_{i\in I}U_i
  =
  X\setminus\bigcap_{i\in I}C_i.
\]
The intersection \(\bigcap_{i\in I}C_i\) belongs to \(\mathcal{C}\), so
\(\bigcup_{i\in I}U_i\in\tau\).  If \(U_1,U_2\in\tau\), write
\(U_1=X\setminus C_1\) and \(U_2=X\setminus C_2\) with
\(C_1,C_2\in\mathcal{C}\).  Then
\[
  U_1\cap U_2
  =
  X\setminus(C_1\cup C_2).
\]
Since \(C_1\cup C_2\in\mathcal{C}\), it follows that \(U_1\cap U_2\in\tau\).
By induction, finite intersections of elements of \(\tau\) belong to \(\tau\).
Thus \(\tau\) is a topology on \(X\).

For this topology, a subset \(F\subseteq X\) is closed exactly when
\(X\setminus F\in\tau\).  By the definition of \(\tau\), this holds exactly
when \(F\in\mathcal{C}\).  Hence the family of closed sets is
\(\mathcal{C}\).

Finally, if \(\sigma\) is any topology on \(X\) whose closed sets are exactly
\(\mathcal{C}\), then its open sets are precisely the complements of elements
of \(\mathcal{C}\).  Therefore
\[
  \sigma=\{X\setminus C:C\in\mathcal{C}\}=\tau,
\]
so the topology is unique.
\end{proof}

\begin{proof}[Detailed Learning Proof]
\LRAProofBodyStart
The closed sets determine the open sets by taking complements.  Therefore the
only possible topology with closed-set family \(\mathcal{C}\) is
\[
  \tau=\{X\setminus C:C\in\mathcal{C}\}.
\]
The assumptions on \(\mathcal{C}\) are exactly the complement-dual versions of
the topology axioms.  The sets \(\emptyset\) and \(X\) are open because they
are complements of \(X\) and \(\emptyset\).  Arbitrary unions of open sets are
open because the complement of an arbitrary union is an arbitrary
intersection.  Finite intersections of open sets are open because the
complement of a finite intersection is a finite union.

Thus \(\tau\) is a topology.  Its closed sets are exactly the sets in
\(\mathcal{C}\), by construction.  Since any topology with closed sets
\(\mathcal{C}\) must have exactly the complements of those sets as its open
sets, no second topology can have the same closed-set family.
\end{proof}

\begin{remark*}[Proof structure]
Construct the only possible topology from complements of the proposed closed
sets.  The three closure assumptions on \(\mathcal{C}\) verify the topology
axioms through De Morgan's laws.  Uniqueness follows because closed sets
determine open sets by complementation.
\end{remark*}

\begin{dependencies}
\begin{itemize}
  \item \hyperref[def:topological-space]{Definition: Topological Space}
  \item \hyperref[def:topological-closed-subset]{Definition: Closed Subset of a Topological Space}
\end{itemize}
\end{dependencies}

\clearpage
